\chapter{Ce que valent les pénalités}\label{ch:penalites}
\iterbanner{V3}

\section{La pénalité avec bits d'écart}
Pour une contrainte $h(\mathbf b)\ge0$ affine en bits, la formulation classique introduit un écart $s\ge0$ et pénalise l'égalité $h=s$ :
\[ P\,\big(h(\mathbf b)-s\big)^2,\qquad s=\frac{s_{\max}}{2^{n_s}-1}\sum_{j=0}^{n_s-1}2^j\,a_j,\quad a_j\in\{0,1\}. \]
Au minimum sur les bits d'écart $a_j$, ce terme vaut $P\,h^2$ si $h<0$ (on paie le carré de la violation) et $P\,(h-s^\star)^2$ si $h\ge0$, où $s^\star$ est le niveau d'écart le plus proche de $h$. Dans notre boucle : $P=75$, 3 bits d'écart, $s_{\max}$ égal à la plus grande marge possible.

\paragraph{Quatre défauts, tous observés.}
\begin{enumerate}
\item \textbf{Une petite violation coûte presque rien.} Une violation de 0,05 sur $C_{L,\max}$ coûte $75\times0{,}05^2\approx0{,}19$, alors que l'objectif normalisé s'étend sur environ 1. Si violer la contrainte fait gagner plus de 0,19 sur l'objectif, l'optimum du QUBO la viole. C'est ce qu'on mesure au chapitre~\ref{ch:hors-ligne} : le candidat retenu a un $\hat C_{L,\max}$ médian de 1,20, \emph{pile} sur la borne, et viole la contrainte du modèle une fois sur deux. Une pénalité quadratique n'est exacte que pour $P\to\infty$ ; une pénalité exacte serait en $|h|$, qui n'est pas quadratique.
\item \textbf{Les dents de scie.} Pour $h\ge0$, l'écart ne prend que 8 valeurs : le terme vaut $P(h-s^\star)^2$, nul seulement aux niveaux de l'écart, jusqu'à $P\Delta^2/4$ entre eux ($\Delta=s_{\max}/7$). Des profils également faisables reçoivent des pénalités différentes, sans raison physique (figure~\ref{fig:penalites}).
\item \textbf{Trois qubits de plus}, couplés à tous les bits de conception : 33 qubits au lieu de 30 à l'échelle, 93 couplages supplémentaires.
\item \textbf{Le spectre est écrasé.} Pour que la pénalité tienne, $P$ doit être grand ; l'objectif n'occupe alors que 1,7~\% (V2) ou 0,8~\% (CST) de l'étendue des énergies. Pour un recuit simulé réglé sur l'objectif, c'est gérable ; pour QAOA ou un recuit quantique, dont la précision est relative à l'étendue totale, c'est rédhibitoire.
\end{enumerate}
Il y a un cinquième défaut, de principe : le carré d'une contrainte quadratique est de degré 4. On ne peut donc pénaliser ainsi qu'une contrainte \emph{affine}, ce qui exclut un modèle quadratique de $C_{L,\max}$ et le classifieur du chapitre~\ref{ch:classifieur}.

\begin{figure}[htbp]
\centering
% figure générée par tools/make_figures.py
\begin{tikzpicture}
\begin{axis}[width=0.92\linewidth,height=6.2cm,xmin=-0.3,xmax=0.5,ymin=-0.6,ymax=2.6,
  xlabel={marge $h=\hat C_{L,\max}-1{,}2$},ylabel={terme ajouté à l'énergie},
  grid=major,grid style={cmuted!20},legend cell align=left,legend style={font=\footnotesize,at={(0.5,-0.3)},anchor=north,legend columns=1},
  extra x ticks={0},extra x tick style={grid style={cink,thick}}]
\fill[corange!8] (axis cs:-0.3,-0.6) rectangle (axis cs:0,2.6);
\node[font=\footnotesize,cink] at (axis cs:-0.15,2.35) {violée};
\node[font=\footnotesize,cink] at (axis cs:0.25,2.35) {respectée};
\addplot[cblue,very thick] coordinates {(-0.3,6.75) (-0.298,6.6603) (-0.296,6.5712) (-0.294,6.4827) (-0.292,6.3948) (-0.29,6.3075) (-0.288,6.2208) (-0.286,6.1347) (-0.284,6.0492) (-0.282,5.9643) (-0.28,5.88) (-0.278,5.7963) (-0.276,5.7132) (-0.274,5.6307) (-0.272,5.5488) (-0.27,5.4675) (-0.268,5.3868) (-0.266,5.3067) (-0.264,5.2272) (-0.262,5.1483) (-0.26,5.07) (-0.258,4.9923) (-0.256,4.9152) (-0.254,4.8387) (-0.252,4.7628) (-0.25,4.6875) (-0.248,4.6128) (-0.246,4.5387) (-0.244,4.4652) (-0.242,4.3923) (-0.24,4.32) (-0.238,4.2483) (-0.236,4.1772) (-0.234,4.1067) (-0.232,4.0368) (-0.23,3.9675) (-0.228,3.8988) (-0.226,3.8307) (-0.224,3.7632) (-0.222,3.6963) (-0.22,3.63) (-0.218,3.5643) (-0.216,3.4992) (-0.214,3.4347) (-0.212,3.3708) (-0.21,3.3075) (-0.208,3.2448) (-0.206,3.1827) (-0.204,3.1212) (-0.202,3.0603) (-0.2,3) (-0.198,2.9403) (-0.196,2.8812) (-0.194,2.8227) (-0.192,2.7648) (-0.19,2.7075) (-0.188,2.6508) (-0.186,2.5947) (-0.184,2.5392) (-0.182,2.4843) (-0.18,2.43) (-0.178,2.3763) (-0.176,2.3232) (-0.174,2.2707) (-0.172,2.2188) (-0.17,2.1675) (-0.168,2.1168) (-0.166,2.0667) (-0.164,2.0172) (-0.162,1.9683) (-0.16,1.92) (-0.158,1.8723) (-0.156,1.8252) (-0.154,1.7787) (-0.152,1.7328) (-0.15,1.6875) (-0.148,1.6428) (-0.146,1.5987) (-0.144,1.5552) (-0.142,1.5123) (-0.14,1.47) (-0.138,1.4283) (-0.136,1.3872) (-0.134,1.3467) (-0.132,1.3068) (-0.13,1.2675) (-0.128,1.2288) (-0.126,1.1907) (-0.124,1.1532) (-0.122,1.1163) (-0.12,1.08) (-0.118,1.0443) (-0.116,1.0092) (-0.114,0.9747) (-0.112,0.9408) (-0.11,0.9075) (-0.108,0.8748) (-0.106,0.8427) (-0.104,0.8112) (-0.102,0.7803) (-0.1,0.75) (-0.098,0.7203) (-0.096,0.6912) (-0.094,0.6627) (-0.092,0.6348) (-0.09,0.6075) (-0.088,0.5808) (-0.086,0.5547) (-0.084,0.5292) (-0.082,0.5043) (-0.08,0.48) (-0.078,0.4563) (-0.076,0.4332) (-0.074,0.4107) (-0.072,0.3888) (-0.07,0.3675) (-0.068,0.3468) (-0.066,0.3267) (-0.064,0.3072) (-0.062,0.2883) (-0.06,0.27) (-0.058,0.2523) (-0.056,0.2352) (-0.054,0.2187) (-0.052,0.2028) (-0.05,0.1875) (-0.048,0.1728) (-0.046,0.1587) (-0.044,0.1452) (-0.042,0.1323) (-0.04,0.12) (-0.038,0.1083) (-0.036,0.0972) (-0.034,0.0867) (-0.032,0.0768) (-0.03,0.0675) (-0.028,0.0588) (-0.026,0.0507) (-0.024,0.0432) (-0.022,0.0363) (-0.02,0.03) (-0.018,0.0243) (-0.016,0.0192) (-0.014,0.0147) (-0.012,0.0108) (-0.01,0.0075) (-0.008,0.0048) (-0.006,0.0027) (-0.004,0.0012) (-0.002,0.0003) (0,0) (0.002,0.0003) (0.004,0.0012) (0.006,0.0027) (0.008,0.0048) (0.01,0.0075) (0.012,0.0108) (0.014,0.0147) (0.016,0.0192) (0.018,0.0243) (0.02,0.03) (0.022,0.0363) (0.024,0.0432) (0.026,0.0507) (0.028,0.0588) (0.03,0.0675) (0.032,0.0768) (0.034,0.0867) (0.036,0.094139) (0.038,0.08381) (0.04,0.074082) (0.042,0.064953) (0.044,0.056424) (0.046,0.048496) (0.048,0.041167) (0.05,0.034439) (0.052,0.02831) (0.054,0.022782) (0.056,0.017853) (0.058,0.013524) (0.06,0.0097959) (0.062,0.0066673) (0.064,0.0041388) (0.066,0.0022102) (0.068,0.00088163) (0.07,0.00015306) (0.072,2.449e-05) (0.074,0.00049592) (0.076,0.0015673) (0.078,0.0032388) (0.08,0.0055102) (0.082,0.0083816) (0.084,0.011853) (0.086,0.015924) (0.088,0.020596) (0.09,0.025867) (0.092,0.031739) (0.094,0.03821) (0.096,0.045282) (0.098,0.052953) (0.1,0.061224) (0.102,0.070096) (0.104,0.079567) (0.106,0.089639) (0.108,0.091127) (0.11,0.080969) (0.112,0.071412) (0.114,0.062455) (0.116,0.054098) (0.118,0.046341) (0.12,0.039184) (0.122,0.032627) (0.124,0.026669) (0.126,0.021312) (0.128,0.016555) (0.13,0.012398) (0.132,0.0088408) (0.134,0.0058837) (0.136,0.0035265) (0.138,0.0017694) (0.14,0.00061224) (0.142,5.5102e-05) (0.144,9.7959e-05) (0.146,0.00074082) (0.148,0.0019837) (0.15,0.0038265) (0.152,0.0062694) (0.154,0.0093122) (0.156,0.012955) (0.158,0.017198) (0.16,0.022041) (0.162,0.027484) (0.164,0.033527) (0.166,0.040169) (0.168,0.047412) (0.17,0.055255) (0.172,0.063698) (0.174,0.072741) (0.176,0.082384) (0.178,0.092627) (0.18,0.088163) (0.182,0.078178) (0.184,0.068792) (0.186,0.060006) (0.188,0.05182) (0.19,0.044235) (0.192,0.037249) (0.194,0.030863) (0.196,0.025078) (0.198,0.019892) (0.2,0.015306) (0.202,0.01132) (0.204,0.0079347) (0.206,0.005149) (0.208,0.0029633) (0.21,0.0013776) (0.212,0.00039184) (0.214,6.1224e-06) (0.216,0.00022041) (0.218,0.0010347) (0.22,0.002449) (0.222,0.0044633) (0.224,0.0070776) (0.226,0.010292) (0.228,0.014106) (0.23,0.01852) (0.232,0.023535) (0.234,0.029149) (0.236,0.035363) (0.238,0.042178) (0.24,0.049592) (0.242,0.057606) (0.244,0.06622) (0.246,0.075435) (0.248,0.085249) (0.25,0.095663) (0.252,0.085249) (0.254,0.075435) (0.256,0.06622) (0.258,0.057606) (0.26,0.049592) (0.262,0.042178) (0.264,0.035363) (0.266,0.029149) (0.268,0.023535) (0.27,0.01852) (0.272,0.014106) (0.274,0.010292) (0.276,0.0070776) (0.278,0.0044633) (0.28,0.002449) (0.282,0.0010347) (0.284,0.00022041) (0.286,6.1224e-06) (0.288,0.00039184) (0.29,0.0013776) (0.292,0.0029633) (0.294,0.005149) (0.296,0.0079347) (0.298,0.01132) (0.3,0.015306) (0.302,0.019892) (0.304,0.025078) (0.306,0.030863) (0.308,0.037249) (0.31,0.044235) (0.312,0.05182) (0.314,0.060006) (0.316,0.068792) (0.318,0.078178) (0.32,0.088163) (0.322,0.092627) (0.324,0.082384) (0.326,0.072741) (0.328,0.063698) (0.33,0.055255) (0.332,0.047412) (0.334,0.040169) (0.336,0.033527) (0.338,0.027484) (0.34,0.022041) (0.342,0.017198) (0.344,0.012955) (0.346,0.0093122) (0.348,0.0062694) (0.35,0.0038265) (0.352,0.0019837) (0.354,0.00074082) (0.356,9.7959e-05) (0.358,5.5102e-05) (0.36,0.00061224) (0.362,0.0017694) (0.364,0.0035265) (0.366,0.0058837) (0.368,0.0088408) (0.37,0.012398) (0.372,0.016555) (0.374,0.021312) (0.376,0.026669) (0.378,0.032627) (0.38,0.039184) (0.382,0.046341) (0.384,0.054098) (0.386,0.062455) (0.388,0.071412) (0.39,0.080969) (0.392,0.091127) (0.394,0.089639) (0.396,0.079567) (0.398,0.070096) (0.4,0.061224) (0.402,0.052953) (0.404,0.045282) (0.406,0.03821) (0.408,0.031739) (0.41,0.025867) (0.412,0.020596) (0.414,0.015924) (0.416,0.011853) (0.418,0.0083816) (0.42,0.0055102) (0.422,0.0032388) (0.424,0.0015673) (0.426,0.00049592) (0.428,2.449e-05) (0.43,0.00015306) (0.432,0.00088163) (0.434,0.0022102) (0.436,0.0041388) (0.438,0.0066673) (0.44,0.0097959) (0.442,0.013524) (0.444,0.017853) (0.446,0.022782) (0.448,0.02831) (0.45,0.034439) (0.452,0.041167) (0.454,0.048496) (0.456,0.056424) (0.458,0.064953) (0.46,0.074082) (0.462,0.08381) (0.464,0.094139) (0.466,0.0867) (0.468,0.0768) (0.47,0.0675) (0.472,0.0588) (0.474,0.0507) (0.476,0.0432) (0.478,0.0363) (0.48,0.03) (0.482,0.0243) (0.484,0.0192) (0.486,0.0147) (0.488,0.0108) (0.49,0.0075) (0.492,0.0048) (0.494,0.0027) (0.496,0.0012) (0.498,0.0003) (0.5,0)};
\addlegendentry{bits d'écart (V2) : $P\min_s(h-s)^2$, 3 bits}
\addplot[corange,very thick] coordinates {(-0.3,9) (-0.298,8.8953) (-0.296,8.7912) (-0.294,8.6877) (-0.292,8.5848) (-0.29,8.4825) (-0.288,8.3808) (-0.286,8.2797) (-0.284,8.1792) (-0.282,8.0793) (-0.28,7.98) (-0.278,7.8813) (-0.276,7.7832) (-0.274,7.6857) (-0.272,7.5888) (-0.27,7.4925) (-0.268,7.3968) (-0.266,7.3017) (-0.264,7.2072) (-0.262,7.1133) (-0.26,7.02) (-0.258,6.9273) (-0.256,6.8352) (-0.254,6.7437) (-0.252,6.6528) (-0.25,6.5625) (-0.248,6.4728) (-0.246,6.3837) (-0.244,6.2952) (-0.242,6.2073) (-0.24,6.12) (-0.238,6.0333) (-0.236,5.9472) (-0.234,5.8617) (-0.232,5.7768) (-0.23,5.6925) (-0.228,5.6088) (-0.226,5.5257) (-0.224,5.4432) (-0.222,5.3613) (-0.22,5.28) (-0.218,5.1993) (-0.216,5.1192) (-0.214,5.0397) (-0.212,4.9608) (-0.21,4.8825) (-0.208,4.8048) (-0.206,4.7277) (-0.204,4.6512) (-0.202,4.5753) (-0.2,4.5) (-0.198,4.4253) (-0.196,4.3512) (-0.194,4.2777) (-0.192,4.2048) (-0.19,4.1325) (-0.188,4.0608) (-0.186,3.9897) (-0.184,3.9192) (-0.182,3.8493) (-0.18,3.78) (-0.178,3.7113) (-0.176,3.6432) (-0.174,3.5757) (-0.172,3.5088) (-0.17,3.4425) (-0.168,3.3768) (-0.166,3.3117) (-0.164,3.2472) (-0.162,3.1833) (-0.16,3.12) (-0.158,3.0573) (-0.156,2.9952) (-0.154,2.9337) (-0.152,2.8728) (-0.15,2.8125) (-0.148,2.7528) (-0.146,2.6937) (-0.144,2.6352) (-0.142,2.5773) (-0.14,2.52) (-0.138,2.4633) (-0.136,2.4072) (-0.134,2.3517) (-0.132,2.2968) (-0.13,2.2425) (-0.128,2.1888) (-0.126,2.1357) (-0.124,2.0832) (-0.122,2.0313) (-0.12,1.98) (-0.118,1.9293) (-0.116,1.8792) (-0.114,1.8297) (-0.112,1.7808) (-0.11,1.7325) (-0.108,1.6848) (-0.106,1.6377) (-0.104,1.5912) (-0.102,1.5453) (-0.1,1.5) (-0.098,1.4553) (-0.096,1.4112) (-0.094,1.3677) (-0.092,1.3248) (-0.09,1.2825) (-0.088,1.2408) (-0.086,1.1997) (-0.084,1.1592) (-0.082,1.1193) (-0.08,1.08) (-0.078,1.0413) (-0.076,1.0032) (-0.074,0.9657) (-0.072,0.9288) (-0.07,0.8925) (-0.068,0.8568) (-0.066,0.8217) (-0.064,0.7872) (-0.062,0.7533) (-0.06,0.72) (-0.058,0.6873) (-0.056,0.6552) (-0.054,0.6237) (-0.052,0.5928) (-0.05,0.5625) (-0.048,0.5328) (-0.046,0.5037) (-0.044,0.4752) (-0.042,0.4473) (-0.04,0.42) (-0.038,0.3933) (-0.036,0.3672) (-0.034,0.3417) (-0.032,0.3168) (-0.03,0.2925) (-0.028,0.2688) (-0.026,0.2457) (-0.024,0.2232) (-0.022,0.2013) (-0.02,0.18) (-0.018,0.1593) (-0.016,0.1392) (-0.014,0.1197) (-0.012,0.1008) (-0.01,0.0825) (-0.008,0.0648) (-0.006,0.0477) (-0.004,0.0312) (-0.002,0.0153) (0,0) (0.002,-0.0147) (0.004,-0.0288) (0.006,-0.0423) (0.008,-0.0552) (0.01,-0.0675) (0.012,-0.0792) (0.014,-0.0903) (0.016,-0.1008) (0.018,-0.1107) (0.02,-0.12) (0.022,-0.1287) (0.024,-0.1368) (0.026,-0.1443) (0.028,-0.1512) (0.03,-0.1575) (0.032,-0.1632) (0.034,-0.1683) (0.036,-0.1728) (0.038,-0.1767) (0.04,-0.18) (0.042,-0.1827) (0.044,-0.1848) (0.046,-0.1863) (0.048,-0.1872) (0.05,-0.1875) (0.052,-0.1872) (0.054,-0.1863) (0.056,-0.1848) (0.058,-0.1827) (0.06,-0.18) (0.062,-0.1767) (0.064,-0.1728) (0.066,-0.1683) (0.068,-0.1632) (0.07,-0.1575) (0.072,-0.1512) (0.074,-0.1443) (0.076,-0.1368) (0.078,-0.1287) (0.08,-0.12) (0.082,-0.1107) (0.084,-0.1008) (0.086,-0.0903) (0.088,-0.0792) (0.09,-0.0675) (0.092,-0.0552) (0.094,-0.0423) (0.096,-0.0288) (0.098,-0.0147) (0.1,2.6021e-16) (0.102,0.0153) (0.104,0.0312) (0.106,0.0477) (0.108,0.0648) (0.11,0.0825) (0.112,0.1008) (0.114,0.1197) (0.116,0.1392) (0.118,0.1593) (0.12,0.18) (0.122,0.2013) (0.124,0.2232) (0.126,0.2457) (0.128,0.2688) (0.13,0.2925) (0.132,0.3168) (0.134,0.3417) (0.136,0.3672) (0.138,0.3933) (0.14,0.42) (0.142,0.4473) (0.144,0.4752) (0.146,0.5037) (0.148,0.5328) (0.15,0.5625) (0.152,0.5928) (0.154,0.6237) (0.156,0.6552) (0.158,0.6873) (0.16,0.72) (0.162,0.7533) (0.164,0.7872) (0.166,0.8217) (0.168,0.8568) (0.17,0.8925) (0.172,0.9288) (0.174,0.9657) (0.176,1.0032) (0.178,1.0413) (0.18,1.08) (0.182,1.1193) (0.184,1.1592) (0.186,1.1997) (0.188,1.2408) (0.19,1.2825) (0.192,1.3248) (0.194,1.3677) (0.196,1.4112) (0.198,1.4553) (0.2,1.5) (0.202,1.5453) (0.204,1.5912) (0.206,1.6377) (0.208,1.6848) (0.21,1.7325) (0.212,1.7808) (0.214,1.8297) (0.216,1.8792) (0.218,1.9293) (0.22,1.98) (0.222,2.0313) (0.224,2.0832) (0.226,2.1357) (0.228,2.1888) (0.23,2.2425) (0.232,2.2968) (0.234,2.3517) (0.236,2.4072) (0.238,2.4633) (0.24,2.52) (0.242,2.5773) (0.244,2.6352) (0.246,2.6937) (0.248,2.7528) (0.25,2.8125) (0.252,2.8728) (0.254,2.9337) (0.256,2.9952) (0.258,3.0573) (0.26,3.12) (0.262,3.1833) (0.264,3.2472) (0.266,3.3117) (0.268,3.3768) (0.27,3.4425) (0.272,3.5088) (0.274,3.5757) (0.276,3.6432) (0.278,3.7113) (0.28,3.78) (0.282,3.8493) (0.284,3.9192) (0.286,3.9897) (0.288,4.0608) (0.29,4.1325) (0.292,4.2048) (0.294,4.2777) (0.296,4.3512) (0.298,4.4253) (0.3,4.5) (0.302,4.5753) (0.304,4.6512) (0.306,4.7277) (0.308,4.8048) (0.31,4.8825) (0.312,4.9608) (0.314,5.0397) (0.316,5.1192) (0.318,5.1993) (0.32,5.28) (0.322,5.3613) (0.324,5.4432) (0.326,5.5257) (0.328,5.6088) (0.33,5.6925) (0.332,5.7768) (0.334,5.8617) (0.336,5.9472) (0.338,6.0333) (0.34,6.12) (0.342,6.2073) (0.344,6.2952) (0.346,6.3837) (0.348,6.4728) (0.35,6.5625) (0.352,6.6528) (0.354,6.7437) (0.356,6.8352) (0.358,6.9273) (0.36,7.02) (0.362,7.1133) (0.364,7.2072) (0.366,7.3017) (0.368,7.3968) (0.37,7.4925) (0.372,7.5888) (0.374,7.6857) (0.376,7.7832) (0.378,7.8813) (0.38,7.98) (0.382,8.0793) (0.384,8.1792) (0.386,8.2797) (0.388,8.3808) (0.39,8.4825) (0.392,8.5848) (0.394,8.6877) (0.396,8.7912) (0.398,8.8953) (0.4,9) (0.402,9.1053) (0.404,9.2112) (0.406,9.3177) (0.408,9.4248) (0.41,9.5325) (0.412,9.6408) (0.414,9.7497) (0.416,9.8592) (0.418,9.9693) (0.42,10.08) (0.422,10.191) (0.424,10.303) (0.426,10.416) (0.428,10.529) (0.43,10.642) (0.432,10.757) (0.434,10.872) (0.436,10.987) (0.438,11.103) (0.44,11.22) (0.442,11.337) (0.444,11.455) (0.446,11.574) (0.448,11.693) (0.45,11.812) (0.452,11.933) (0.454,12.054) (0.456,12.175) (0.458,12.297) (0.46,12.42) (0.462,12.543) (0.464,12.667) (0.466,12.792) (0.468,12.917) (0.47,13.043) (0.472,13.169) (0.474,13.296) (0.476,13.423) (0.478,13.551) (0.48,13.68) (0.482,13.809) (0.484,13.939) (0.486,14.07) (0.488,14.201) (0.49,14.333) (0.492,14.465) (0.494,14.598) (0.496,14.731) (0.498,14.865) (0.5,15)};
\addlegendentry{déséquilibrée : $P(h^2-2h^\star h)$, $h^\star=0{,}05$}
\addplot[cgreen,very thick] coordinates {(-0.3,0.24) (-0.298,0.2384) (-0.296,0.2368) (-0.294,0.2352) (-0.292,0.2336) (-0.29,0.232) (-0.288,0.2304) (-0.286,0.2288) (-0.284,0.2272) (-0.282,0.2256) (-0.28,0.224) (-0.278,0.2224) (-0.276,0.2208) (-0.274,0.2192) (-0.272,0.2176) (-0.27,0.216) (-0.268,0.2144) (-0.266,0.2128) (-0.264,0.2112) (-0.262,0.2096) (-0.26,0.208) (-0.258,0.2064) (-0.256,0.2048) (-0.254,0.2032) (-0.252,0.2016) (-0.25,0.2) (-0.248,0.1984) (-0.246,0.1968) (-0.244,0.1952) (-0.242,0.1936) (-0.24,0.192) (-0.238,0.1904) (-0.236,0.1888) (-0.234,0.1872) (-0.232,0.1856) (-0.23,0.184) (-0.228,0.1824) (-0.226,0.1808) (-0.224,0.1792) (-0.222,0.1776) (-0.22,0.176) (-0.218,0.1744) (-0.216,0.1728) (-0.214,0.1712) (-0.212,0.1696) (-0.21,0.168) (-0.208,0.1664) (-0.206,0.1648) (-0.204,0.1632) (-0.202,0.1616) (-0.2,0.16) (-0.198,0.1584) (-0.196,0.1568) (-0.194,0.1552) (-0.192,0.1536) (-0.19,0.152) (-0.188,0.1504) (-0.186,0.1488) (-0.184,0.1472) (-0.182,0.1456) (-0.18,0.144) (-0.178,0.1424) (-0.176,0.1408) (-0.174,0.1392) (-0.172,0.1376) (-0.17,0.136) (-0.168,0.1344) (-0.166,0.1328) (-0.164,0.1312) (-0.162,0.1296) (-0.16,0.128) (-0.158,0.1264) (-0.156,0.1248) (-0.154,0.1232) (-0.152,0.1216) (-0.15,0.12) (-0.148,0.1184) (-0.146,0.1168) (-0.144,0.1152) (-0.142,0.1136) (-0.14,0.112) (-0.138,0.1104) (-0.136,0.1088) (-0.134,0.1072) (-0.132,0.1056) (-0.13,0.104) (-0.128,0.1024) (-0.126,0.1008) (-0.124,0.0992) (-0.122,0.0976) (-0.12,0.096) (-0.118,0.0944) (-0.116,0.0928) (-0.114,0.0912) (-0.112,0.0896) (-0.11,0.088) (-0.108,0.0864) (-0.106,0.0848) (-0.104,0.0832) (-0.102,0.0816) (-0.1,0.08) (-0.098,0.0784) (-0.096,0.0768) (-0.094,0.0752) (-0.092,0.0736) (-0.09,0.072) (-0.088,0.0704) (-0.086,0.0688) (-0.084,0.0672) (-0.082,0.0656) (-0.08,0.064) (-0.078,0.0624) (-0.076,0.0608) (-0.074,0.0592) (-0.072,0.0576) (-0.07,0.056) (-0.068,0.0544) (-0.066,0.0528) (-0.064,0.0512) (-0.062,0.0496) (-0.06,0.048) (-0.058,0.0464) (-0.056,0.0448) (-0.054,0.0432) (-0.052,0.0416) (-0.05,0.04) (-0.048,0.0384) (-0.046,0.0368) (-0.044,0.0352) (-0.042,0.0336) (-0.04,0.032) (-0.038,0.0304) (-0.036,0.0288) (-0.034,0.0272) (-0.032,0.0256) (-0.03,0.024) (-0.028,0.0224) (-0.026,0.0208) (-0.024,0.0192) (-0.022,0.0176) (-0.02,0.016) (-0.018,0.0144) (-0.016,0.0128) (-0.014,0.0112) (-0.012,0.0096) (-0.01,0.008) (-0.008,0.0064) (-0.006,0.0048) (-0.004,0.0032) (-0.002,0.0016) (0,-0) (0.002,-0.0016) (0.004,-0.0032) (0.006,-0.0048) (0.008,-0.0064) (0.01,-0.008) (0.012,-0.0096) (0.014,-0.0112) (0.016,-0.0128) (0.018,-0.0144) (0.02,-0.016) (0.022,-0.0176) (0.024,-0.0192) (0.026,-0.0208) (0.028,-0.0224) (0.03,-0.024) (0.032,-0.0256) (0.034,-0.0272) (0.036,-0.0288) (0.038,-0.0304) (0.04,-0.032) (0.042,-0.0336) (0.044,-0.0352) (0.046,-0.0368) (0.048,-0.0384) (0.05,-0.04) (0.052,-0.0416) (0.054,-0.0432) (0.056,-0.0448) (0.058,-0.0464) (0.06,-0.048) (0.062,-0.0496) (0.064,-0.0512) (0.066,-0.0528) (0.068,-0.0544) (0.07,-0.056) (0.072,-0.0576) (0.074,-0.0592) (0.076,-0.0608) (0.078,-0.0624) (0.08,-0.064) (0.082,-0.0656) (0.084,-0.0672) (0.086,-0.0688) (0.088,-0.0704) (0.09,-0.072) (0.092,-0.0736) (0.094,-0.0752) (0.096,-0.0768) (0.098,-0.0784) (0.1,-0.08) (0.102,-0.0816) (0.104,-0.0832) (0.106,-0.0848) (0.108,-0.0864) (0.11,-0.088) (0.112,-0.0896) (0.114,-0.0912) (0.116,-0.0928) (0.118,-0.0944) (0.12,-0.096) (0.122,-0.0976) (0.124,-0.0992) (0.126,-0.1008) (0.128,-0.1024) (0.13,-0.104) (0.132,-0.1056) (0.134,-0.1072) (0.136,-0.1088) (0.138,-0.1104) (0.14,-0.112) (0.142,-0.1136) (0.144,-0.1152) (0.146,-0.1168) (0.148,-0.1184) (0.15,-0.12) (0.152,-0.1216) (0.154,-0.1232) (0.156,-0.1248) (0.158,-0.1264) (0.16,-0.128) (0.162,-0.1296) (0.164,-0.1312) (0.166,-0.1328) (0.168,-0.1344) (0.17,-0.136) (0.172,-0.1376) (0.174,-0.1392) (0.176,-0.1408) (0.178,-0.1424) (0.18,-0.144) (0.182,-0.1456) (0.184,-0.1472) (0.186,-0.1488) (0.188,-0.1504) (0.19,-0.152) (0.192,-0.1536) (0.194,-0.1552) (0.196,-0.1568) (0.198,-0.1584) (0.2,-0.16) (0.202,-0.1616) (0.204,-0.1632) (0.206,-0.1648) (0.208,-0.1664) (0.21,-0.168) (0.212,-0.1696) (0.214,-0.1712) (0.216,-0.1728) (0.218,-0.1744) (0.22,-0.176) (0.222,-0.1776) (0.224,-0.1792) (0.226,-0.1808) (0.228,-0.1824) (0.23,-0.184) (0.232,-0.1856) (0.234,-0.1872) (0.236,-0.1888) (0.238,-0.1904) (0.24,-0.192) (0.242,-0.1936) (0.244,-0.1952) (0.246,-0.1968) (0.248,-0.1984) (0.25,-0.2) (0.252,-0.2016) (0.254,-0.2032) (0.256,-0.2048) (0.258,-0.2064) (0.26,-0.208) (0.262,-0.2096) (0.264,-0.2112) (0.266,-0.2128) (0.268,-0.2144) (0.27,-0.216) (0.272,-0.2176) (0.274,-0.2192) (0.276,-0.2208) (0.278,-0.2224) (0.28,-0.224) (0.282,-0.2256) (0.284,-0.2272) (0.286,-0.2288) (0.288,-0.2304) (0.29,-0.232) (0.292,-0.2336) (0.294,-0.2352) (0.296,-0.2368) (0.298,-0.2384) (0.3,-0.24) (0.302,-0.2416) (0.304,-0.2432) (0.306,-0.2448) (0.308,-0.2464) (0.31,-0.248) (0.312,-0.2496) (0.314,-0.2512) (0.316,-0.2528) (0.318,-0.2544) (0.32,-0.256) (0.322,-0.2576) (0.324,-0.2592) (0.326,-0.2608) (0.328,-0.2624) (0.33,-0.264) (0.332,-0.2656) (0.334,-0.2672) (0.336,-0.2688) (0.338,-0.2704) (0.34,-0.272) (0.342,-0.2736) (0.344,-0.2752) (0.346,-0.2768) (0.348,-0.2784) (0.35,-0.28) (0.352,-0.2816) (0.354,-0.2832) (0.356,-0.2848) (0.358,-0.2864) (0.36,-0.288) (0.362,-0.2896) (0.364,-0.2912) (0.366,-0.2928) (0.368,-0.2944) (0.37,-0.296) (0.372,-0.2976) (0.374,-0.2992) (0.376,-0.3008) (0.378,-0.3024) (0.38,-0.304) (0.382,-0.3056) (0.384,-0.3072) (0.386,-0.3088) (0.388,-0.3104) (0.39,-0.312) (0.392,-0.3136) (0.394,-0.3152) (0.396,-0.3168) (0.398,-0.3184) (0.4,-0.32) (0.402,-0.3216) (0.404,-0.3232) (0.406,-0.3248) (0.408,-0.3264) (0.41,-0.328) (0.412,-0.3296) (0.414,-0.3312) (0.416,-0.3328) (0.418,-0.3344) (0.42,-0.336) (0.422,-0.3376) (0.424,-0.3392) (0.426,-0.3408) (0.428,-0.3424) (0.43,-0.344) (0.432,-0.3456) (0.434,-0.3472) (0.436,-0.3488) (0.438,-0.3504) (0.44,-0.352) (0.442,-0.3536) (0.444,-0.3552) (0.446,-0.3568) (0.448,-0.3584) (0.45,-0.36) (0.452,-0.3616) (0.454,-0.3632) (0.456,-0.3648) (0.458,-0.3664) (0.46,-0.368) (0.462,-0.3696) (0.464,-0.3712) (0.466,-0.3728) (0.468,-0.3744) (0.47,-0.376) (0.472,-0.3776) (0.474,-0.3792) (0.476,-0.3808) (0.478,-0.3824) (0.48,-0.384) (0.482,-0.3856) (0.484,-0.3872) (0.486,-0.3888) (0.488,-0.3904) (0.49,-0.392) (0.492,-0.3936) (0.494,-0.3952) (0.496,-0.3968) (0.498,-0.3984) (0.5,-0.4)};
\addlegendentry{lagrangien : $-\lambda h/s$, $\lambda/s=0{,}8$}
\end{axis}
\end{tikzpicture}

\caption{Les trois termes ajoutés à l'énergie, en fonction de la marge $h$ ($P=75$, échelle de $C_{L,\max}$). Bits d'écart : violation payée au carré, dents de scie du côté faisable. Pénalité déséquilibrée : pas de qubit, mais minimum en $h^\star$ et croissance du côté faisable. Lagrangien : une droite, qui récompense la marge au lieu de la pénaliser.}\label{fig:penalites}
\end{figure}

\section{La pénalité déséquilibrée}
Montañez-Barrera, Willsch, Maldonado-Romo et Michielsen (\emph{Unbalanced penalization}, Quantum Science and Technology, 2024) suppriment les bits d'écart. L'idée part de la pénalité idéale $e^{-h}$, grande pour $h<0$ et presque nulle pour $h>0$, et en garde le développement de degré 2 : $1-h+h^2/2$. Le terme ajouté est $-\lambda_1h+\lambda_2h^2$, quadratique si $h$ est affine. Nous l'écrivons
\[ P\,\big(h^2-2h^\star h\big)=P\,\big[(h-h^\star)^2-h^{\star2}\big], \]
nulle en $h=0$, minimale en $h=h^\star>0$, positive pour $h<0$. Elle a été testée sur le voyageur de commerce, le sac à dos et le rangement en boîtes ; sur ce dernier, elle permet de traiter sur machine D-Wave des instances de 29 objets, contre 11 avec les bits d'écart.

\paragraph{Son défaut chez nous.} Elle pénalise aussi le côté faisable : avec $P=75$ et $h^\star=0{,}05$, un profil de marge $0{,}3$ paie $75\times(0{,}09-0{,}03)=4{,}5$, quatre fois l'étendue de l'objectif. Elle attire donc tous les candidats vers $h\approx h^\star$, ce qui déforme l'objectif (les « faux minimums » relevés dans la littérature). Sur des problèmes combinatoires où la contrainte est saturée à l'optimum (un sac à dos plein), c'est sans conséquence ; sur un front de Pareto dont une partie est loin de la contrainte, c'est un biais.

\section{Régler les poids : ce que dit la littérature}
Le réglage de $P$ est un sujet de recherche à lui seul. Les méthodes publiées se rangent en trois familles :
\begin{itemize}
\item \textbf{bornes exactes} : un $P$ assez grand pour qu'aucune violation ne soit rentable, calculé à partir des coefficients de l'objectif. Sûr, mais il écrase le spectre ;
\item \textbf{poids séquentiels} : partir d'un petit $P$ et l'augmenter tant que les échantillons sont infaisables (travaux de l'université d'Exeter, d'Ayodele) ;
\item \textbf{poids par contrainte} réglés sur le taux de faisabilité des échantillons, par exemple avec une loi sigmoïde.
\end{itemize}
Notre $P=75$, fixé à la main sur le problème V2, n'appartient à aucune. La deuxième famille est très proche du lagrangien du chapitre suivant, qui ajuste lui aussi un poids d'après les échantillons, mais sur un terme linéaire.

\begin{resultat}
\begin{itemize}
\item La pénalité quadratique avec bits d'écart rend les petites violations rentables, crée des dents de scie, ajoute 3 qubits et écrase le spectre. Elle ne s'applique qu'à une contrainte affine.
\item La pénalité déséquilibrée supprime les bits d'écart, mais attire tous les candidats vers une marge fixée : elle déforme l'objectif loin de la contrainte.
\item Toute pénalité au carré est limitée aux contraintes affines.
\end{itemize}
\end{resultat}

\exercices
\begin{exercice}\label{ex:pen-rentable}
Objectif normalisé : un profil faisable de marge $h=0$ vaut $-0{,}80$ ; son voisin vaut $-1{,}00$ avec $h=-0{,}04$. (a) Avec la pénalité à bits d'écart ($P=75$), lequel est l'optimum du QUBO ? (b) Quel $P$ faudrait-il pour que le profil faisable gagne ? (c) Que devient alors la part de l'objectif dans le spectre, si la marge prédite varie de $-0{,}5$ à $+0{,}5$ sur la grille ?
\end{exercice}
\begin{exercice}\label{ex:pen-deseq}
Pénalité déséquilibrée avec $P=75$, $h^\star=0{,}05$. (a) Calculer le terme en $h=-0{,}05$, $h=0$, $h=0{,}05$, $h=0{,}3$. (b) Pour quelles marges $h>0$ la pénalité est-elle positive ? (c) Proposer $P$ et $h^\star$ pour que le terme reste inférieur à 0,5 jusqu'à $h=0{,}3$, et dire ce que l'on perd.
\end{exercice}
\begin{exercice}\label{ex:pen-degre}
Le modèle de $C_{L,\max}$ devient quadratique en $z$. Montrer sur l'exemple $h=z_1z_2-0{,}2$ (deux variables à 1 bit, $z_i=2b_i-1$) que $h^2$ contient un monôme de degré supérieur à 2 en bits, sauf simplification due à $b_i^2=b_i$. Conclure pour des variables à 3 bits.
\end{exercice}

\begin{pistes}
\begin{itemize}
\item Régler $P$ par la méthode séquentielle (augmenter tant que la meilleure lecture viole la contrainte) au lieu de $P=75$ : c'est la comparaison la plus juste entre bits d'écart et lagrangien.
\end{itemize}
\end{pistes}
